\subsection{DECOMPOSITION OF A LIE ALGEBRA RELATIVE TO AN AUTOMORPHISM}

PROPOSITION 12. Let \(\mathfrak{g}\) be a Lie algebra, a an automorphism of \(\mathfrak{g}\) .

\begin{enumerate}
\def\labelenumi{(\roman{enumi})}
\item
  For \(\lambda, \mu \in k\) , \([\mathfrak{g}^{\lambda}(a), \mathfrak{g}^{\mu}(a)] \subset \mathfrak{g}^{\lambda \mu}(a)\) ; in particular, \(\mathfrak{g}^1(a)\) is a subalgebra of \(\mathfrak{g}\) .
\item
  If B is a symmetric bilinear form on g invariant under a, \(\mathfrak{g}^{\lambda}(a)\) and \(\mathfrak{g}^{\mu}(a)\) are orthogonal relative to B for \(\lambda\mu \neq 1\) . Assume that B is nondegenerate. Then, if \(\lambda \neq 0\) , the restriction of B to \(\mathfrak{g}^{\lambda}(a) \times \mathfrak{g}^{1/\lambda}(a)\) is nondegenerate.
\end{enumerate}

Assertion (i) and the first half of (ii) follow from Prop. 2 (iii) applied to the composition law \(\mathfrak{g} \times \mathfrak{g} \to \mathfrak{g}\) and the bilinear form B. To prove the second half of (ii), we can assume that \(k\) is algebraically closed. Then \(\mathfrak{g} = \bigoplus_{\nu \in k} \mathfrak{g}^{\nu}(a)\) . In view of the above, \(\mathfrak{g}^{\lambda}(a)\) is orthogonal to \(\mathfrak{g}^{\nu}(a)\) if \(\lambda \nu \neq 1\) ; since B is non-degenerate, it follows that its restriction to \(\mathfrak{g}^{\lambda}(a) \times \mathfrak{g}^{1/\lambda}(a)\) is also.

COROLLARY. Assume that \(k\) is of characteristic zero and that \(\mathfrak{g}\) is semi-simple. Then the subalgebra \(\mathfrak{g}^1(a)\) satisfies conditions (1) and (2) of Lemma 2; it is reductive in \(\mathfrak{g}\) .

Condition (1) follows from part (ii) of Prop. 12; condition (2) follows from Prop. 4 of no. 1.

